%==============================================================================
\chapter{Constraints and technical requirements}
%==============================================================================

\section{Architecture}
A \textbf{scalable, evolvable, flexible, multi-tier and loosely coupled}
architecture, relying on \textbf{Spring Boot 3} (Java 17 LTS) as the main
infrastructure, with a strict separation between the service API and the
presentation client.

\section{Technical requirements adopted}
\begin{xltabular}{\textwidth}{>{\bfseries}p{4cm} X}
\toprule
\textbf{Requirement} & \textbf{Description} \\
\midrule
\endhead
Layered architecture & Controller / service / persistence separation, dependency inversion provided by the Spring container. \\
REST API & Spring MVC controllers exposing a versioned API, described by an \textbf{OpenAPI 3} contract. \\
Presentation & \textbf{React} client in \textbf{TypeScript}, decoupled from the server and consuming the REST API. \\
Business services & Spring components (\texttt{@Service}) encapsulating the business rules. \\
Persistence & \textbf{Spring Data JPA} for CRUD, \textbf{jOOQ} for analytical and geospatial queries. \\
Internationalisation & UI texts externalised in per-locale resource files, open to offline translation with no language limit, FR, EN and AR at minimum in the demonstration. \\
Technical configuration & Parameterisation via Spring profiles (\texttt{application.yml}) and environment variables: data sources, endpoints, secrets. \\
Business configuration & Thresholds, units and active modules in the text file \texttt{ConfigZumm.ini}, adjustable by the operator without recompilation or redeployment. \\
Inter-component security & \textbf{X.509} certificates and \textbf{TLS 1.3} protocol. \\
User authentication & \textbf{OpenID Connect} via \textbf{Keycloak}, with Google identity federation and \textbf{JWT} tokens. \\
Third-party API  & REST endpoint exposing the \texttt{getZummHoneyActualQuantity(int zummID)} operation to retrieve the honey quantity of a hive, published in the OpenAPI contract. \\
Front-end        & JavaScript / TypeScript and free libraries, respecting the imposed usability. \\
\bottomrule
\end{xltabular}

\begin{encadre}
\textbf{Identifier stability:} the change of technical foundation alters neither
the business domain nor the public identifiers. The
\texttt{getZummHoneyActualQuantity} operation keeps its name and signature,
whatever transport technology is chosen to expose it.
\end{encadre}

\section{Non-functional requirements}
\begin{itemize}
  \item \textbf{Usability, user-friendliness and simplicity} of use, interface consistency and quality of the graphic design (of paramount importance).
  \item \textbf{Scalability}: openness to adding languages, units and indicators.
  \item \textbf{Interoperability}: API queryable by third-party applications from a published contract.
  \item \textbf{Security}: encryption of exchanges and strong authentication delegated to an identity provider.
  \item \textbf{Design quality}: compliance with the \textbf{SOLID} principles and use of \textbf{design patterns} to guarantee loose coupling and scalability, as detailed in appendix~\ref{chap:solid}.
\end{itemize}

\begin{encadre}
\textbf{Design principles adopted:} the architecture applies the five
\textbf{SOLID} principles (single responsibility, open/closed, Liskov
substitution, interface segregation, dependency inversion) and a set of
\textbf{design patterns} (Composite, State, Strategy, Observer, Command, Factory
Method, Builder, Adapter, Facade, Singleton, Proxy). Their application to the
system, justified ``before / after'', is detailed in
appendix~\ref{chap:solid}.
\end{encadre}

\section{Design questions to cover}
The design must explicitly answer the following questions:
\begin{enumerate}
  \item What are the \textbf{inputs} (list, data dictionary and types) and \textbf{outputs} (list, dictionary, types, formulas and calculation queries) of the system?
  \item What does the solution consist of?
  \item Who are the typical users and their access rights to the features?
  \item How does the system carry out the user operations?
  \item How is the system packaged and deployed?
  \item How do the elements interact and in what order / chronology?
  \item Why, how and where was each technology used?
\end{enumerate}

\section{Justification of the technology choices}
The following table answers the why, the how and the where of each adopted
technology. The \textbf{concrete implementations} of this foundation (runtime,
DBMS, front-end libraries, API style) and the \textbf{justified comparison} of
the evaluated alternatives are detailed in
appendix~\ref{chap:technologies}.
\begin{xltabular}{\textwidth}{>{\bfseries}p{2.6cm} X p{5.4cm}}
\toprule
\textbf{Technology} & \textbf{Why} & \textbf{Where} \\
\midrule
\endhead
Spring Boot 3 & Market-standard application foundation, injection container and autoconfiguration & General organisation of the application \\
Spring MVC & Process the HTTP requests and orchestrate the processing & Control layer \\
React + TypeScript & Produce a rich, typed and testable interface & Presentation layer \\
Spring Data JPA & Access the data without hand-writing the CRUD & Persistence layer \\
jOOQ & Express in typed SQL the analytical and PostGIS queries that JPA renders poorly & Dashboards and geoprocessing \\
OpenAPI 3 & Publish a machine-readable contract and generate third-party clients & Interoperability interface \\
i18n resources & Externalise the texts for offline translation & Internationalisation module \\
Spring profiles & Parameterise the infrastructure per environment without rebuilding the image & Technical configuration \\
ConfigZumm.ini & Let the operator adjust business thresholds without technical intervention & Business configuration \\
Keycloak / OIDC & Authenticate and authorise without managing passwords or re-coding RBAC & Access to the system \\
X.509 and TLS 1.3 & Encrypt and authenticate the exchanges & Inter-component communications \\
\bottomrule
\end{xltabular}

\section{Packaging and deployment}
\begin{itemize}
  \item \textbf{Packaging}: the server application is packaged as an \textbf{executable Java archive} (self-contained JAR, embedded Tomcat server), then as a \textbf{container image} including the internationalisation resource files. The React client is built into static files served by the reverse proxy.
  \item \textbf{Runtime}: deployment of orchestrated containers, with no external application server to administer.
  \item \textbf{Database}: \textbf{PostgreSQL} extended by \textbf{PostGIS} (geolocation) and \textbf{TimescaleDB} (time series of measurements).
  \item \textbf{Security}: \textbf{TLS} termination at the reverse proxy with X.509 certificates, no application or management port published directly, configuration of the Keycloak identity realm.
  \item \textbf{Configuration}: the infrastructure is set through Spring profiles and environment variables; the business thresholds, units and active modules remain adjusted in \texttt{ConfigZumm.ini}, mounted as a volume and reloadable without redeploying the image.
\end{itemize}
The physical distribution is illustrated by the deployment diagram (figure \ref{fig:deploiement}).
